% ============================================================
% Come si usano, argomento 1: eventi e insiemi
% ============================================================
\argomento{Eventi e insiemi}{Formule a p.~\pageref*{f:operazioni}}

\begin{scelta}{Il primo passo di ogni esercizio: tradurre il testo in insiemi}
Si dà un nome agli eventi semplici ($A$, $B$, \ldots) e si traduce la frase del testo in operazioni fra insiemi. La
tabella è il vocabolario.
\begin{center}
\renewcommand{\arraystretch}{1.25}
\begin{tabular}{@{}ll@{}}
  \toprule
  \textbf{Nel testo} & \textbf{Con gli insiemi} \\
  \midrule
  «$A$ oppure $B$», «almeno uno dei due» & $A \cup B$ \\
  «$A$ e $B$», «entrambi» & $A \cap B$ \\
  «non $A$» & $\compl{A}$ \\
  «$A$ ma non $B$», «solo $A$ (fra i due)» & $A \setminus B = A \cap \compl{B}$ \\
  «né $A$ né $B$», «nessuno dei due» & $\compl{A} \cap \compl{B} = \compl{A \cup B}$ \\
  «non entrambi», «al più uno dei due» & $\compl{A} \cup \compl{B} = \compl{A \cap B}$ \\
  «$A$ e $B$ non possono verificarsi insieme» & $A \cap B = \emptyset$ \\
  «se si verifica $A$, si verifica anche $B$» & $A \subseteq B$ \\
  «succede sempre» / «non succede mai» & $\Omega$ / $\emptyset$ \\
  \bottomrule
\end{tabular}
\end{center}
\end{scelta}

\scheda{operazioni}{
  \usoquando{Sempre, all'inizio: per scrivere in simboli l'evento di cui il testo chiede la probabilità.}
  \usocome{Si individuano gli eventi semplici e si collegano con il vocabolario: «o» diventa $\cup$, «e» diventa
    $\cap$, «non» diventa il complementare.}
  \usoattenzione{La «o» della matematica è inclusiva: $A \cup B$ comprende anche i risultati in cui si verificano
    entrambi.}
  \usoesempio{Un dado: «esce un pari oppure almeno 5» è $\{2, 4, 6\} \cup \{5, 6\} = \{2, 4, 5, 6\}$; il 6 si
    prende una volta sola.}
}

\scheda{differenza}{
  \usoquando{Il testo dice «$A$ ma non $B$»; oppure serve trasformare una differenza in un'intersezione, per
    applicarle le regole delle intersezioni (De Morgan, indipendenza).}
  \usoesempio{$A_1 = \{1, \ldots, 8\}$, $A_2 = \{1, 4, 6\}$: \ $A_1 \setminus A_2 = A_1 \cap \compl{A_2} =
    \{2, 3, 5, 7, 8\}$.}
}

\scheda{complementare}{
  \usoquando{Per semplificare un'espressione, e ogni volta che si divide il mondo in due casi: $A$ e $\compl{A}$
    coprono tutto $\Omega$ senza sovrapporsi. È la partizione più semplice (\vv{partizione}).}
  \usoesempio{Un bit trasmesso arriva giusto oppure sbagliato, e non può essere le due cose insieme:
    $G \cup \compl{G} = \Omega$, $G \cap \compl{G} = \emptyset$.}
}

\scheda{demorgan}{
  \usoquando{Il testo dice «né\ldots{} né\ldots», «nessuno», «non entrambi»; oppure si deve passare dal
    complementare di un'unione a un'intersezione, o viceversa.}
  \usocome{Si fa il complementare di ogni evento e si scambia $\cup$ con $\cap$. Serve spesso con eventi
    indipendenti, perché la probabilità di un'intersezione è un prodotto (\vv{indipendenza}) e quella di
    un'unione no.}
  \usovarianti{Con $M$ eventi la regola è la stessa: «nessuno degli $A_i$» è $\bigcap_i \compl{A_i}$.}
  \usoattenzione{$\compl{A \cup B}$ non è $\compl{A} \cup \compl{B}$: insieme ai complementari va cambiato anche il
    simbolo.}
  \usoesempio{Due bit, $A_i$ = «il bit $i$ è sbagliato». «Almeno uno sbagliato» è $A_1 \cup A_2$; il suo
    complementare, «nessuno sbagliato», è $\compl{A_1 \cup A_2} = \compl{A_1} \cap \compl{A_2}$.}
}

\scheda{associativa}{
  \usoquando{Per togliere le parentesi: si scrive $A \cup B \cup C$ («almeno uno dei tre») e $A \cap B \cap C$
    («tutti e tre»).}
}

\scheda{distributiva}{
  \usoquando{Per spezzare un evento in pezzi che non si sovrappongono, a cui poi si applica l'additività:
    $A = A \cap (B \cup \compl{B}) = (A \cap B) \cup (A \cap \compl{B})$. È il passaggio chiave della probabilità
    totale (\vv{totale}).}
  \usoattenzione{La distributiva di $\cup$ rispetto a $\cap$ non ha l'analoga fra i numeri:
    $a + b\,c \neq (a + b)(a + c)$. Con gli insiemi $\cup$ e $\cap$ si comportano in modo simmetrico.}
}

\scheda{cartesiano}{
  \usoquando{Per scrivere lo spazio campione di un esperimento fatto di più prove: due lanci di una moneta danno
    $\{T, C\}^2$, tre lanci di un dado $\{1, \ldots, 6\}^3$.}
  \usoattenzione{Le coppie sono ordinate: $(T, C)$ e $(C, T)$ sono due risultati diversi. Contarli come uno solo
    porta a probabilità sbagliate (\vv{laplace}).}
  \usoesempio{Due lanci di un dado: $\card{\Omega} = 6 \cdot 6 = 36$ coppie ordinate, tutte equiprobabili.}
}
